\section{Evaluacion heuristica y depuracion del pipeline visual}
\label{sec:evaluacion-heuristica-depuracion}

\subsection{Ausencia de \textit{ground truth} anotado}

Durante el desarrollo no se dispuso de un conjunto de datos anotado frame a frame que pudiera usarse como \textit{ground truth} para evaluar de forma supervisada todo el pipeline de procesamiento visual. Por ese motivo, la validacion del sistema se apoyo en una estrategia heuristica, combinando inspeccion visual y metricas agregadas tras cada ejecucion.

En particular, la inspeccion cualitativa se realizo sobre un video de prueba de referencia\footnote{\url{https://github.com/hectorgarciaa/narrador-futbol/blob/main/data/partidoPrueba/partido.mp4}}, revisando la evolucion temporal de las identidades y la consistencia de los tracks en situaciones dificiles (oclusiones, cruces, perdidas temporales y reapariciones).

\subsection{Metricas heuristicas principales}

Tras ejecutar el tracking, el sistema calcula un resumen de metricas para jugadores y balon. Entre ellas se incluyen la cobertura media, que mide durante que proporcion del video aparece cada ID; el numero medio de fragmentos, que indica cuantas interrupciones tiene una trayectoria; la longitud media de los huecos sin deteccion; la velocidad media y maxima estimada a partir del movimiento de las \textit{bounding boxes}; la tasa de cambios de equipo asignado; la entropia del equipo, que refleja la inestabilidad en la asignacion; la variacion del color de camiseta; la variacion del tamano de la \textit{bounding box}; y la confianza media de las detecciones asociadas a cada track.

\subsection{Uso de metricas para ajuste de hiperparametros}

Estas metricas fueron especialmente utiles para guiar la depuracion de todos los modulos visuales y no solo del tracking canonico. Por ejemplo, velocidades maximas excesivamente altas solian indicar intercambios de identidad entre jugadores lejanos; una cobertura baja podia senalar perdidas frecuentes de deteccion o \textit{gates} demasiado estrictos; y un numero alto de fragmentos apuntaba a trayectorias poco continuas.

Del mismo modo, cambios frecuentes de equipo o una alta variacion del color de camiseta podian revelar errores de clasificacion o reasignacion. Por tanto, la evaluacion del pipeline visual se apoyo en dos fuentes complementarias: inspeccion visual del video generado y analisis cuantitativo de metricas. La primera permitia entender el tipo de error y su contexto, mientras que la segunda permitia comparar ejecuciones y verificar si una modificacion mejoraba realmente la estabilidad global del sistema.

